\documentclass[12pt,a4paper,UTF8]{ctexart}
\usepackage[a4paper,top=22mm,bottom=22mm,left=22mm,right=22mm,headheight=15pt]{geometry}
\newcommand{\Needspace}[1]{\par\begingroup\dimen0=#1\relax\dimen2=\pagegoal\advance\dimen2 by-\pagetotal\ifdim\dimen0>\dimen2\newpage\fi\endgroup}
\newcommand{\needspace}[1]{\par\begingroup\skip0=\lastskip\vskip-\skip0\dimen0=#1\relax\advance\dimen0 by\skip0\vskip\dimen0\penalty-100\vskip-\dimen0\vskip\skip0\penalty9999\endgroup}
\usepackage{amsmath,amssymb,mathtools,bm}
\usepackage{graphicx,xcolor,booktabs,longtable,array,enumitem,fancyhdr,caption,float,microtype,setspace}
\usepackage[hidelinks,unicode]{hyperref}
\usepackage[most]{tcolorbox}
\definecolor{ContestBlue}{HTML}{17365D}
\hypersetup{pdftitle={第057套},pdfauthor={}}
\setstretch{1.24}
\setlength{\parindent}{2em}
\setlength{\parskip}{0.35em}
\setlength{\emergencystretch}{2em}
\setlist{itemsep=0.35em,topsep=0.4em}
\captionsetup{font=small,labelfont=bf,justification=centering}
\setcounter{secnumdepth}{0}
\pagestyle{fancy}
\fancyhf{}
\fancyhead[L]{\small 第057套}
\fancyfoot[C]{\thepage}
\renewcommand{\headrulewidth}{0.4pt}
\ctexset{
  section={format=\Large\bfseries\color{ContestBlue},beforeskip=1.2em,afterskip=0.7em},
  subsection={format=\large\bfseries\color{ContestBlue},beforeskip=1.0em,afterskip=0.5em},
  subsubsection={format=\normalsize\bfseries,beforeskip=0.8em,afterskip=0.35em}
}
\usepackage{calc}
\newcounter{none}
\let\originalsubsection\subsection
\renewcommand{\subsection}{\needspace{6\baselineskip}\originalsubsection}
\makeatletter
\newcommand{\tightlist}{\setlength{\itemsep}{0pt}\setlength{\parskip}{0pt}}
\newcommand{\maxwidth}{\ifdim\Gin@nat@width>\linewidth\linewidth\else\Gin@nat@width\fi}
\newcommand{\maxheight}{\ifdim\Gin@nat@height>0.72\textheight 0.72\textheight\else\Gin@nat@height\fi}
\newcommand{\pandocbounded}[1]{\begingroup\setkeys{Gin}{width=\maxwidth,height=\maxheight,keepaspectratio}#1\endgroup}
\makeatother
\begin{document}
\begin{titlepage}
  \centering
  \vspace*{58mm}
  {\Huge\bfseries\color{ContestBlue} 第057套\par}
  \vfill
  \vspace*{28mm}
\end{titlepage}
\section{第 1 题}\label{ux7b2c-1-ux9898}

如图所示，质量分别为 \(m_A\)、\(m_B\) 的小球 \(A\)、\(B\) 由长度为 \(L\)
的轻直杆连接。小球 \(A\) 被限制在内壁光滑的竖直滑槽中运动，小球 \(B\)
位于光滑水平面上。初始时杆竖直，两球均静止。受微小扰动后，\(B\)
开始沿水平面向右运动。求：

（1）\(A\) 球着地时的速度；

（2）\(A\) 球机械能最小时，水平面对 \(B\) 球的支持力；

（3）若 \(m_A=m_B\)，求 \(A\) 球机械能最小时杆与竖直方向夹角 \(\theta\)
的余弦；

（4）在第（3）问条件下，求 \(A\)
球机械能的最小值，以水平面为重力势能零面。

\begin{figure}[H]
\centering
\includegraphics[width=0.68\linewidth,height=\textheight,keepaspectratio,alt={图1-1　轻杆连接的两球及其相互垂直的运动约束}]{../figures/question-01.pdf}
\caption*{图1-1　轻杆连接的两球及其相互垂直的运动约束}
\end{figure}

\newpage

\section{第 2 题}\label{ux7b2c-2-ux9898}

天花板上有两个定滑轮，间距为 \(l\)。一根总长为 \(4l\)
的轻绳依次绕过两个定滑轮，并在中间两段上各悬挂一个质量为 \(m\)
的动滑轮。所有滑轮均可视为质点，忽略摩擦。初始时两个动滑轮均静止在与两个定滑轮构成正三角形的平衡位置
\(A\)，且二者位置重合。绳始终绷紧，运动保持在初始平面内。求：

（1）把其中一个动滑轮固定在 \(A\)，另一动滑轮在 \(A\) 附近小振动的频率；

（2）限制两个动滑轮都只能沿通过 \(A\) 的竖直直线运动，求它们在 \(A\)
附近上下小振动的频率。

\begin{figure}[H]
\centering
\includegraphics[width=0.76\linewidth,height=\textheight,keepaspectratio,alt={图2-1　两定滑轮与两动滑轮构成的闭合绳路}]{../figures/question-02.pdf}
\caption*{图2-1　两定滑轮与两动滑轮构成的闭合绳路}
\end{figure}

\newpage

\section{第 3 题}\label{ux7b2c-3-ux9898}

在水平面内有一根质量为 \(m\)、长度为 \(2l\)
的光滑刚性细杆，可无摩擦地绕穿过其中点 \(O\)
的竖直轴转动。空间中存在与转轴平行、方向垂直纸面向里的匀强磁场
\(\bm B\)。把质量为 \(2m\)、电量为 \(q>0\)
的光滑小环套在杆上。小环与杆之间无电荷转移，不计相对论效应。求：

（1）杆连同小环以角速度 \(\omega_0\)
转动时，小环在杆上任意位置均可相对杆静止，求无量纲比值
\(\omega_0m/(qB)\)；

（2）杆的初角速度为
\(\omega_1\)，小环初始静止在转轴处。受微小扰动后小环沿杆运动，且恰好不从杆端脱离，求
\(\omega_1\)；

（3）初角速度为 \(\omega_2\)
时，小环从杆端脱出。能否使小环在杆转过半圈之前重新回到杆上并如此往复？令
\(\alpha=\omega_2/(qB/m)\)，列出判定所用的数值方程，并说明结论。数值量保留四位有效数字。

\begin{figure}[H]
\centering
\includegraphics[width=0.74\linewidth,height=\textheight,keepaspectratio,alt={图3-1　匀强磁场中的转动细杆与带电小环}]{../figures/question-03.pdf}
\caption*{图3-1　匀强磁场中的转动细杆与带电小环}
\end{figure}

\newpage

\section{第 4 题}\label{ux7b2c-4-ux9898}

研究图示交流电桥。毫伏表按输入阻抗足够大的电压表处理。忽略导线电阻、接触电阻及一切互感。左侧待测线圈的电感为
\(L_x\)、内阻为 \(R_L\)；右侧标准线圈的电感为 \(L_0\)、内阻为
\(R_{L0}\)；其余三只电阻的阻值为 \(R_0,R_1,R_2\)。

（1）接通交流电源后，若交流毫伏表读数为零，且除 \(L_x,R_L\)
外其余量均已知，用已知量表示 \(L_x,R_L\)。

（2）若毫伏表读数不为零，且 \(R_{L0}\) 固定，\(R_0,R_1,R_2,L_0\)
均可连续调节，证明至少需要选择两个独立调节量才能一般性地调平。实际器材中，电源角频率为
\(\omega=10^4\ \mathrm{s^{-1}}\)，待测线圈标称
\(L_x=\sqrt3\ \mathrm{mH}\)、\(R_L=30\ \Omega\)，\(R_1=\sqrt3\ \mathrm{k}\Omega\)，初始
\(R_0=10\ \Omega\)、\(R_2=1.00\ \mathrm{k}\Omega\)；标准线圈调至
\(L_0=2.00\ \mathrm{mH}\)，此时
\(R_{L0}=(20/\sqrt3-10)\ \Omega\)。每一组操作先只调 \(R_0\)
使读数最小，再只调 \(R_2\)
使读数最小，如此反复。至少经过多少组操作，\(L_x\)
测量值与标称值的相对误差才能小于 \(0.03\)？

\begin{figure}[H]
\centering
\includegraphics[width=0.613\linewidth,height=\textheight,keepaspectratio,alt={图4-1　含待测线圈与标准线圈的交流电桥}]{../figures/question-04.pdf}
\caption*{图4-1　含待测线圈与标准线圈的交流电桥}
\end{figure}

\newpage

\section{第 5 题}\label{ux7b2c-5-ux9898}

一个圆形平行板电容器的极板面积为 \(S\)，初始板间距为 \(d\)，且 \(d\)
远小于圆盘半径。每块极板质量均为
\(m\)，各由中间一根光滑绝缘直杆导向，只能沿竖直方向运动。两根导杆始终与相应极板保持良好电接触，并由可伸缩导线连接到恒压电源
\(U\)。两极板的总电荷为 \(2Q\)
且保持不变；电源在两板间转移电荷以维持电势差
\(U\)。装置处于真空中，真空介电常数为 \(\varepsilon_0\)、光速为
\(c\)。忽略边缘效应、重力以及导杆对电场和磁场的影响。

（1）释放后若两极板能在静电作用下保持静止，求 \(Q\)；

（2）判断该平衡是否稳定；

（3）以下取第（1）问的 \(Q\)，在 \(t=0\)
时给一微扰，使两极板开始相互远离。任意时刻板距为
\(x(t)\)，求两板相对速度 \(v=\dot x\)；

（4）求板间电场强度随时间的变化率 \(\dot E(x)\)，保留正负号；

（5）在两圆盘正中间的平面上，距导杆为 \(r\) 的点满足 \(r\)
远小于极板半径。求该点磁感应强度的大小 \(B(x)\)；

（6）一个面积为 \(A\)、均匀带电 \(+q\)
的绝缘小圆环，其圆心位于两极板的中间平面，圆心到导杆的距离为
\(r\)。环面包含圆心处的径向方向和极板法线方向，环法线沿环向磁场方向；环的尺寸远小于
\(r\) 及板间距。小环相对其轴的转动惯量为
\(I\)，只可绕自身轴转动，忽略小环自感。初始角速度为零。求小环可达到的最大角速度
\(\omega_0\) 以及此时板距 \(d_0\)。

\begin{figure}[H]
\centering
\includegraphics[width=0.72\linewidth,height=\textheight,keepaspectratio,alt={图5-1　恒压连接的平行圆盘电容器与竖直导向结构}]{../figures/question-05.pdf}
\caption*{图5-1　恒压连接的平行圆盘电容器与竖直导向结构}
\end{figure}

\newpage

\section{第 6 题}\label{ux7b2c-6-ux9898}

一条理想高分子链由 \(N\) 个 Kuhn 链段组成，每段长度为
\(b\)。链首、链尾距离为 \(R\)，并满足

\[
\langle R^2\rangle=Nb^2.
\]

令末端位置矢量为\(\boldsymbol R\)、其长度为R。高斯链末端位置的三维概率密度为

\[
p(R)=\frac{\Omega(N,R)}
{\int\Omega(N,\boldsymbol R)\,\mathrm d^3R}
=\left(
\frac{3}{2\pi\langle R^2\rangle}
\right)^{3/2}
\exp\left[
-\frac{3R^2}{2\langle R^2\rangle}
\right].
\]

亥姆霍兹自由能为 \(F=U-TS\)，其中 \(U\) 与 \(R\)
无关，\(S=k\ln\Omega\)。求：

（1）相对 \(R=0\) 的初态，证明熵变 \(\Delta S\propto R^2\)；

（2）等温条件下分子链的弹力 \(f=\partial F/\partial R\)；

（3）若 \(N\) 不变而温度变为原来的两倍，作用力如何变化？若温度不变而
\(N\) 变为原来的三倍，作用力如何变化？

\begin{figure}[H]
\centering
\includegraphics[width=0.76\linewidth,height=\textheight,keepaspectratio,alt={图6-1　高斯链的末端距离与链段构象}]{../figures/question-06.pdf}
\caption*{图6-1　高斯链的末端距离与链段构象}
\end{figure}

\newpage

\section{第 7 题}\label{ux7b2c-7-ux9898}

考虑一团低速、纯暗物质组成的稀薄云团，可视为只受自身万有引力作用的理想气体。云团球对称分布并处于力学平衡和热平衡，温度恒为
\(T\)，摩尔质量为 \(\mu\)。距中心 \(r\) 处的质量密度为
\(\rho(r)\)、压强为 \(p(r)\)、引力场大小为 \(g(r)\)，万有引力常量为
\(G\)。

（1）利用理想气体状态方程，用 \(p,\mu,R,T\) 表示 \(\rho\)；

（2）用 \(g,p,\mu,r,T,R\) 表示 \(\mathrm dp/\mathrm dr\)；

（3）建立 \(\rho(r)\) 满足的常微分方程；

（4）质量为 \(m\) 的小星体只受云团引力，在半径 \(r_0\)
处作圆周运动。已知该处压强为 \(p_0\)、引力场为 \(g_0\)，求轨道周期
\(T_0\)；

（5）在圆周运动基础上施加小的径向扰动，求径向振动周期 \(T_1\)。

\begin{figure}[H]
\centering
\includegraphics[width=0.66\linewidth,height=\textheight,keepaspectratio,alt={图7-1　球对称等温自引力气体云团}]{../figures/question-07.pdf}
\caption*{图7-1　球对称等温自引力气体云团}
\end{figure}

\newpage

\section{第 8 题}\label{ux7b2c-8-ux9898}

在惯性系 \(S\) 中，半径为 \(r\) 的圆环上均匀分布 \(N\)
个电子，每个电子电量为 \(-e\)，且 \(N\) 很大。电子以角速度 \(\omega\)
逆时针作圆周运动，满足 \(\omega r<c\)。圆环中心有一个静止的正电荷
\(+Ne\)，因而系统总电荷为零。圆环上点 \(A\) 的位置矢量与 \(x\) 轴夹角为
\(\theta\)。惯性系 \(S'\) 相对 \(S\) 以速度 \(\beta c\) 沿 \(+x\)
方向运动。求：

（1）在 \(S'\) 中，\(A\) 点电子的速度；

（2）在 \(S'\) 中，\(A\) 点圆环电子的线密度；

（3）在 \(S'\) 中，整个系统的磁偶极矩和电偶极矩。

\begin{figure}[H]
\centering
\includegraphics[width=0.64\linewidth,height=\textheight,keepaspectratio,alt={图8-1　中性旋转电子圆环及沿 x 方向的洛伦兹变换}]{../figures/question-08.pdf}
\caption*{图8-1　中性旋转电子圆环及沿 \(x\) 方向的洛伦兹变换}
\end{figure}
\end{document}
